\section{江南朝廷的继承、边防与文化网络}
\label{sec:southern-courts-liu-song-qi-liang}

四五〇年刘宋文帝发动元嘉北伐，宋军自淮、泗北进，太武帝随后南下反击，战线又收缩至淮、泗一带。行军路线、兵数和杀伤存在史书夸饰，北伐受挫与沿江防御被重新动员则是可靠的后果。对建康而言，北方并非遥远的战场：淮北城戍、江防、将领任用和都城戒严都要消耗同一批人力与运输资源。

\subsection{宗室内战如何改变州镇}

四五三年太子刘劭弑杀文帝，刘骏自江州起兵，攻入建康后即位；次年又平定刘义宣、臧质等叛乱。人物动机与每一次结盟不能由胜者纪传一概裁定，不过江州、荆州和建康的军府在继承时能各自调动资源，是这一系列内战的必要条件。孝武帝此后加强对宗室王镇与州郡兵权的约束，却不能使地方军镇从此失去作用。

四六四年至四六七年，幼帝、近侍、宗室与明帝之间再次发生政变和子勋之乱。北魏乘刘宋内战夺取淮北若干城镇，南方州郡则重新编置守将，流民与逃户问题更加突出。把这些变化写成某位皇帝“暴虐”或“英明”的报应，无法解释为何战争会改变户籍、道路和地方武装的分布。宫廷需要禁军和近臣，州镇需要粮饷与名义授权，两者在连续继承危机中越来越难互信。

\subsection{刘宋终局与南齐的建国}

四七五年明帝死后，萧道成掌握禁军；四七七年废后废帝、立顺帝，反对者相继失败，四七九年受禅建南齐。受禅及其前后的任命、赦令和封赏可以由《宋书》《南齐书》《资治通鉴》相互核对，个人是否早有唯一的篡夺计划则不是材料能够直接回答的问题。新朝首先要接收的是宋末州镇、降附军队与文书网络，而不是一片已经服从的江南。

齐武帝在位时以礼制、儒学、选举和州镇任命巩固建康，也要维持淮南防线并同北魏使节往来。建康士族、文人和州郡官僚的活动可从传记和诗文看见，不能被当作全体江南社会的缩影。陶弘景辞官入句曲山、陆修静整理道经的活动，则提示士族仕宦、道教仪式和医药知识能够在朝廷之外形成持续网络；文本的成书、抄写和所叙事态须分别考察。

\subsection{梁的信誉与淮汉防线}

五〇一年萧衍自襄阳起兵，次年受禅建梁。宽刑、整饬吏治和整编齐末降军的初政，是重建建康信誉的方式，也须依赖漕运、户籍与州郡文书的日常接续。梁武帝扶持寺院、讲经和僧官相关活动，使佛教成为礼制之外的重要公共网络；僧祐编《弘明集》《出三藏记集》，保存辩难、经录和僧传的文本传统。它们说明佛教的知识和组织在江南可见，却不能据此断言全社会以同一方式参与。

北魏以寿阳为基地经营淮南屯戍，五〇五年至五〇六年进攻义阳、汉水方向，梁军在钟离据淮水防御。胜负和大致年月可以校核，兵额与战场细节不宜据传世叙事写实。钟离之后的城戍、水运和守将调整，显示南朝的边防是交通网络而非一条静止的长江屏障。到五二〇年代，寺院、城市斋会与同泰寺的舍身叙事又让皇帝、僧团与建康居民的关系更紧密也更复杂；后世禅宗传记关于菩提达摩的若干场景，尤其不能当作当时的逐字记录。

\subsection{史料的视野}

刘宋、南齐和梁的帝纪、列传主要见于《宋书》《南齐书》《梁书》及《南史》，同北魏材料和《资治通鉴》交叉可辨认战争与继承的顺序；它们皆偏向都城、宗室与显贵。道教经目、佛教经录、墓志和城市考古能补正知识、宗教和地方网络，却不等于普通家庭的完整档案。本节因此不以宫廷道德或文人趣味概括江南，而将继承、州镇、边防、宗教与文书视作同一套持续调整的条件。
